\chapter{Simulation Design and Data}
\label{ch:simulation}

\section{Overall Research Design}
\label{sec:design}

The thesis uses a two-part mixed-method design. Part~I,
\Cref{ch:design}, is a conceptual and architectural analysis answering
RQ1 by theoretical synthesis supported by literature. Part~II, this
chapter and the next, is a quantitative simulation study answering RQ2.

The two parts are linked in both directions. The design blocks derived
in Part~I define the parameters Part~II's engines instantiate:
\Cref{sec:interval}'s batch-interval argument and
\Cref{sec:rationing}'s rationing-rule taxonomy are both operationalised
into one fixed configuration, described in \Cref{sec:environment} and
run in \Cref{ch:results} --- not, as a fuller version of this design
would allow, into the interval sweep and rationing-rule comparison that
\Cref{sec:experiments} describes as future work rather than as executed
here. In the other direction, the measured results feed into the
architectural recommendation of \Cref{ch:discussion}, and the
predictions of \Cref{tab:predictions} that the simulation cannot test
define the boundary of what that recommendation can claim.

The identifying idea of Part~II is simple. A single historical order
flow is replayed twice, once through a continuous engine and once
through a batch engine. Everything that could otherwise confound a
comparison between mechanisms --- the asset, the participant
population, the information environment, the time period, the
volatility regime --- is held fixed by construction, because it is
literally the same event stream. The matching engine is the only
treatment variable.

\section{Simulation Environment}
\label{sec:environment}

\subsection{Shared Type System and Matching Interface}

Both engines are implemented in Rust and share one order, trade, and
book type system, and both implement one common matching interface.
This is not a stylistic choice. If the two engines had separate order
representations or separate event loops, any measured difference could
originate in the representation rather than in the mechanism. A shared
interface makes the engine substitutable at a single point in the
harness and makes the claim ``the engine is the only difference''
verifiable rather than asserted.

\begin{lstlisting}[caption={The shared matching interface. Both engines
implement the same trait, so the harness is generic over the
mechanism.},label={lst:trait}]
pub trait MatchingEngine {
    /// Submit, modify or cancel. Returns events generated
    /// immediately (empty for the FBA during accumulation).
    fn on_event(&mut self, ev: &InboundEvent) -> Vec<EngineEvent>;

    /// Advance the engine clock. The FBA clears here when a
    /// batch boundary is crossed; the CDA is a no-op.
    fn on_clock(&mut self, now: Nanos) -> Vec<EngineEvent>;

    /// Snapshot of the current book state for metric computation.
    fn book_state(&self, pair: PairId) -> BookState;

    /// Engine identity and configuration, written to the log
    /// header for reproducibility.
    fn config(&self) -> EngineConfig;
}
\end{lstlisting}

\subsection{The Continuous Engine (CDA)}

The continuous engine maintains a price-time priority limit order book
per asset pair. Each incoming order is matched immediately against the
opposite side of the book, executing at the resting maker's own limit
price, so a single aggressive order that sweeps several price levels can
print several different trade prices in one submission. A limit order's
unfilled remainder rests on the book, ordered within its price level by
arrival; a market order's unfilled remainder is discarded rather than
resting. Time-in-force is parsed from the replayed data but not enforced
beyond this: there is no separate reject-if-crossing path for post-only
orders and no cancel-remainder-on-partial-fill path for immediate-or-cancel
orders. The engine implements two operations, submission and
cancellation; there is no distinct order-modify operation, so this is not
a design choice made to match a specific venue convention but a scope
limitation of the current implementation, stated here rather than left
implicit.

This engine is the control condition. It is not intended to be a novel
contribution but a reference implementation, and \Cref{sec:validation}
describes the invariant checks and hand-computed test cases used to
verify its correctness.

\subsection{The Batch Engine (FBA)}

The batch engine buffers all events arriving within an interval $\tau$
and clears at the end of the interval. Clearing proceeds in the stages
of \Cref{fig:pipeline}: the accumulated buy and sell schedules are
constructed as in \eqref{eq:demandsupply} from the submitted limit
prices only (a market order contributes to demand or supply at every
candidate price but is never itself a candidate price; a batch containing
only market orders therefore has no candidate price and clears nothing,
rather than inventing one), the volume-maximising price is found by a
search over these candidates as in \eqref{eq:clearing}, and ties are
resolved by a fixed four-step cascade: maximise matched volume, then
minimise the residual imbalance, then minimise distance from the
\emph{auction's own previous clearing price} --- the third step of
\Cref{sec:clearing}'s cascade, using the first of its two named
alternatives, the previous clearing price, not an external feed, since
none is available in this dataset (\Cref{sec:data}) --- and finally,
if still tied, the lowest price, for a deterministic result. Inframarginal
orders are filled completely and the marginal side is rationed
price-time-within-the-batch: eligible orders are ranked by aggressiveness
(a market order ranks ahead of any limit order), then by submission time,
then by order id, and filled in that order until the matched volume is
exhausted. This is one instance of the rationing rules \Cref{sec:rationing}
describes, not a configurable choice in the current implementation.

Volume that cannot be matched at $\clr$ is not executed; it always rolls
over into the next batch, since the implementation has no
cancel-on-no-clear path. No external liquidity source, such as an
automated market maker, is used to fill the residual. This is what keeps
the comparison strictly like-for-like: both engines can only execute
against the order flow that was actually submitted, so neither has access
to liquidity the other lacks.

\Cref{ch:design} derives a wider design space for the batch
engine --- the interval length, the rationing rule, the tie-break
cascade, and the boundary policy are all, in principle, separate levers.
The implementation used in this thesis fixes all of them to one
configuration, summarised in \Cref{app:config}: the mechanism just
described, at $\tau = 1$ second. \Cref{sec:experiments} states this
plainly rather than presenting it as one point in a sweep.

\subsection{The Replay Harness}

A shared harness reads the event stream, dispatches each event to the
engine under test in timestamp order, triggers clock ticks so that the
FBA's batch boundaries fall where the configuration says they should,
and writes every order, trade, and book state transition to a structured
log for offline metric computation. Metrics are never computed inside
the engine; they are computed from the log by a separate analysis
pipeline, so that both engines are measured by exactly the same code.

Two properties of the harness matter for the validity of the results.
It is deterministic: given the same input stream and the same
configuration, including the seed for any randomised rule, it produces
byte-identical logs. And it separates simulation time from wall-clock
time, so that engine throughput and clearing latency can be measured
without those measurements affecting the simulated market.

\section{Data Source and Data Levels}
\label{sec:data}

\subsection{The Hyperliquid Level~4 Dataset}

The simulation is fed with the publicly released Hyperliquid Level~4
order book dataset of \citet{albers2026openbook}, captured directly at a
non-validating Hyperliquid node and published on Zenodo. Beyond a
standard L4 feed it also records rejected orders, failed cancellations,
both counterparties to every trade, and post-trade inventory, all at
nanosecond precision. Because it is captured at the node rather than
sampled from periodic snapshots, no book reconstruction step is
required, which removes a substantial source of error that affects
studies working from vendor feeds.

Three properties make this dataset the right choice here. Order-by-order
granularity is the necessary condition for replaying historical flow
through a different matching engine at all. Participant attribution is
what makes the allocation and per-participant metrics of RQ2.3
computable. And nanosecond timestamps are what make interval boundary
placement meaningful at the millisecond scale of the batch intervals
tested.

\subsection{The Data-Level Framework}

\Cref{tab:levels} sets out the standard hierarchy of order book data and
what each level enables. The framework serves two purposes in this
thesis. It scopes the metric catalogue, since every metric in
\Cref{sec:metrics} is tagged with the minimum level it requires. And it
makes the contribution of the dataset explicit by showing which metrics
would simply be unavailable at a lower level.

\begin{table}[htbp]
\centering
\caption[Order book data levels]{The order book data-level hierarchy
and what each level enables. This thesis works with L4 throughout.}
\label{tab:levels}
\small
\begin{tabularx}{\textwidth}{@{}p{1.0cm} p{4.2cm} L@{}}
\toprule
\textbf{Level} & \textbf{Content} & \textbf{What it enables}\\
\midrule
L1 & Best bid and offer, last trade & Quoted spread, midpoint series,
realized volatility, effective spread\\
L2 & Aggregated depth per price level (the ladder) & Depth at best and
within $x$ bps, depth curve shape, book imbalance, slippage of a
hypothetical order\\
L3 & Order-by-order feed: every add, modify, cancel and execution with
an order ID and timestamp & Exact book reconstruction, queue position,
time priority effects, cancellation and order-to-trade ratios,
arrival-time distribution inside an interval\\
L4 & L3 plus participant attribution (wallet or account identifier) per
order and per fill & Per-participant metrics: maker inventory and
markout PnL, adverse selection by counterparty type, fill distribution
across participants under a rationing rule, order size inflation, flow
classification into benign and toxic\\
\bottomrule
\end{tabularx}
\end{table}

\subsection{What the Simulation Observes and What It Proxies}

Because the engines are the author's own, several quantities that are
unobservable in field data become directly observable here: every
submitted limit price including those of orders that never execute, the
exact rationing outcome order by order, the internal clearing state
including the full demand and supply schedules at every candidate price,
and the volume left unexecuted after each batch. Field studies must
infer these or do without them.

What remains unobservable is proxied or modelled instead, and stating
this explicitly is part of the design:

\begin{itemize}
  \item \textbf{The true fundamental value} has no proxy in this dataset.
  The order-status archive used here carries no oracle or mark-price
  feed, so the pricing-error metric of \Cref{tab:metrics2} is defined
  but structurally uncomputable: every value is missing by construction,
  not merely unmeasured. \Cref{app:formulas} states this explicitly
  where the metric is defined, rather than leaving an empty column
  unexplained.
  \item \textbf{Hidden liquidity} that was never submitted as a visible
  order cannot be recovered and is absent from both engines equally.
  \item \textbf{True participant latency} is not observable; the
  timestamps record when the node saw the message, not when the
  participant decided to send it.
\end{itemize}

The third of these bears on interpretation. Because arrival times are
those produced under Hyperliquid's own continuous engine, the position
of an order relative to a simulated batch boundary is an artefact of
where the boundary was placed, not a choice the participant made. The
consequences are drawn out in \Cref{sec:testable}.

\section{Experimental Design and Comparability Protocol}
\label{sec:experiments}

\subsection{Treatment and Controls}

The treatment is the matching engine. The comparability protocol
enforces that as little else as possible differs, and states plainly
where full control was not achieved:

\begin{enumerate}
  \item \textbf{Identical input.} Both runs consume the same event
  stream, byte for byte, in the same order: the full December~2025 SOL
  order-status archive of \Cref{sec:data}, from its first record onward.
  \item \textbf{Identical initial state.} Both engines start from an
  empty book at the same first timestamp. There is no lead-in period
  excluded from metric computation; every bucket from the first record
  onward, including the first, is in scope. A warm-up exclusion was
  considered but not implemented, and is not needed for a
  month-long run.
  \item \textbf{Identical metric code.} Metrics are computed offline
  from the logged CSV output by one implementation (\Cref{sec:metrics},
  \Cref{app:formulas}) shared by both engines.
  \item \textbf{An uncontrolled edge case, stated rather than hidden.}
  Neither engine implements self-trade prevention: an order that would
  fill against the same participant's own resting order executes like
  any other match. This is a genuine limitation of the current
  implementation, not a policy the two engines were made to share on
  purpose, and it is listed here so that it is not mistaken for a
  controlled condition.
  \item \textbf{Determinism.} A single uninterrupted run reproduces
  byte-identical logs from the same input. The run behind
  \Cref{ch:results} was checkpointed and resumed several times over its
  multi-day wall-clock duration (\Cref{app:repro}); a resumed run can
  leave a short gap of empty interval rows at the seam between the last
  file completed before interruption and the first file processed after
  resumption, which is a documented property of the checkpoint format,
  not a source of nondeterminism in any individual computed row.
\end{enumerate}

\subsection{Sample}

The sample is a single instrument, SOL/USD, replayed once through each
engine over the full calendar month the dataset covers: 1--31~December
2025, 744 hourly files, no gaps (\Cref{sec:data}). This is a census of
the available month for this instrument, not a sample drawn along
instrument, session, or time-of-day dimensions --- there is exactly one
of each, so between-instrument and between-session comparisons that a
multi-instrument design would support are not available here.
\Cref{tab:sample} reports its composition.

\subsection{Scope of the Executed Comparison}

The design space \Cref{ch:design} derives has several free dimensions:
the interval $\tau$, the rationing rule, the tie-break cascade, and the
boundary policy. Exploring that space --- an interval sweep across
$\tau$, and a comparison across rationing rules at fixed $\tau$ --- is
the natural next experiment, and \Cref{sec:futurework} returns to it.
It was not run for this thesis. What was run is a single configuration,
fixed for the reasons of scope and compute cost given there: $\tau = 1$
second, the fixed boundary and price-time-within-batch rationing that
\Cref{sec:environment} describes, roll-over residual, and the four-step
tie-break cascade anchored to the auction's own previous clearing price.
Every result in \Cref{ch:results} is this one configuration; none of it
should be read as a point on a sweep.

\subsection{Statistical Treatment}

The comparison is paired by construction: each one-second interval of
the replayed month produces one CDA observation and one FBA observation
from the same flow. Differences are therefore evaluated as paired
differences per interval rather than as two independent samples, which
removes the common variation driven by the market state and leaves the
mechanism effect. Because microstructure series are strongly
autocorrelated, inference uses Newey--West standard errors robust to
autocorrelation, with the lag length chosen by the standard automatic
bandwidth rule and capped for bounded computation
(\Cref{app:formulas}). With a single instrument and a single
continuous session there is no per-instrument or per-session breakdown
to report; every paired-difference result in \Cref{ch:results} is
pooled over the full month, and only over the subset of intervals in
which the metric in question was computable on both sides (most metrics
here are conditional on a trade or a priced batch occurring in the
interval, so this is typically a minority of the 2,678,400 one-second
buckets in the month --- \Cref{ch:results} reports how many for each
metric).

\section{Metric Catalogue}
\label{sec:metrics}

Notation: $a_t$ and $b_t$ are the best ask and bid, $\midp = (a_t+b_t)/2$
the midpoint, $p_k$ the execution price of trade $k$, $q_k$ its size,
$D_k \in \{+1,-1\}$ its sign (buyer- or seller-initiated), $\pi_k$ the
submitted limit price, $\clr$ the uniform clearing price of a batch, and
$\Delta$ a markout horizon. The \emph{Data} column gives the minimum
data level required. \Cref{app:formulas} gives the full, implementation-level
formula for every metric below, including bucketing conventions and the
exact `None`-versus-zero convention each column follows, plus a small
number of further diagnostic metrics computed alongside these but not
tabulated here, notably Kyle's lambda, used in \Cref{ch:results}.

\begin{table}[htbp]
\centering
\caption[Metrics for RQ2.1]{RQ2.1 --- liquidity and transaction cost
metrics.}
\label{tab:metrics1}
\small
\begin{tabularx}{\textwidth}{@{}p{3.1cm} L p{1.6cm}@{}}
\toprule
\textbf{Metric} & \textbf{Definition and measurement} & \textbf{Data}\\
\midrule
Quoted spread & $(a_t-b_t)/\midp$, time-weighted per interval; for the
FBA the implied spread between the marginal unexecuted buy and sell of
the batch & L1/L2\\
Depth at best & Displayed volume at the best bid and ask, time-weighted;
for the FBA the volume at $\clr$ in the batch schedules & L2\\
Depth within $x$ bps & Cumulative volume within $x$ basis points of
$\midp$ (or of $\clr$), for several $x$ & L2\\
Book imbalance & $(V^{bid}-V^{ask})/(V^{bid}+V^{ask})$ at the top $n$
levels & L2\\
Effective spread & $2D_k(p_k-m_k)/m_k$, volume-weighted; $m_k$ is the
pre-trade midpoint, for the FBA the midpoint at the start of the batch
& L1 + trades\\
Realized spread & $2D_k(p_k-m_{k+\Delta})/m_k$ for
$\Delta \in \{1,5,30\}$ seconds; the part of the spread the liquidity
provider keeps & L1 + trades\\
Price impact & Effective spread minus realized spread; the adverse
selection component & L1 + trades\\
Amihud illiquidity & $|r_t|/\text{volume}_t$ averaged over intervals &
L1 + trades\\
\bottomrule
\end{tabularx}
\end{table}

\begin{table}[htbp]
\centering
\caption[Metrics for RQ2.2]{RQ2.2 --- price discovery and market quality
metrics.}
\label{tab:metrics2}
\small
\begin{tabularx}{\textwidth}{@{}p{3.1cm} L p{1.6cm}@{}}
\toprule
\textbf{Metric} & \textbf{Definition and measurement} & \textbf{Data}\\
\midrule
Realized volatility & Standard deviation of interval midpoint or
clearing price returns & L1\\
Intra-interval price dispersion & Standard deviation of execution prices
within one interval; zero by construction for the FBA, strictly positive
for the CDA, and therefore a direct measure of the uniform-price
property & L1 + trades\\
Pricing error & $|\clr-p^{ref}|$ or $|\midp-p^{ref}|$ against an
external reference price, in bps; \emph{not computable in this
dataset} --- see \Cref{sec:data}, \Cref{app:formulas} & L1 + ref\\
\bottomrule
\end{tabularx}
\end{table}

\begin{table}[htbp]
\centering
\caption[Metrics for RQ2.3]{RQ2.3 --- execution, allocation and
order-arrival timing metrics, plus engine controls.}
\label{tab:metrics3}
\small
\begin{tabularx}{\textwidth}{@{}p{3.1cm} L p{1.6cm}@{}}
\toprule
\textbf{Metric} & \textbf{Definition and measurement} & \textbf{Data}\\
\midrule
Executed volume & Matched quantity and notional per interval; number of
trades & L1 + trades\\
Fill rate & Filled quantity divided by submitted quantity, overall and
by participant class & L3/L4\\
Time to execution & Elapsed time from submission to fill; for the FBA
bounded below by the time to the next batch boundary, which quantifies
the delay cost of batching & L3\\
Trader surplus & $\sum_k |\pi_k - \clr|\, q_k$ over executed orders: the
gain relative to the submitted limit price & L3/L4\\
Order size inflation & Ratio of submitted to filled size per participant
over time; the empirical signature of order stuffing under pro rata &
L4\\
Order-to-trade ratio & Submitted messages per executed trade;
cancellation rate per interval; measures message traffic and quote
flickering & L3\\
Boundary concentration & Share of order arrivals in the final $x\%$ of
the batch interval & L3\\
Throughput & Orders processed per second under identical load &
control\\
Clearing latency & Wall-clock time per batch clearing computation and
per continuous match & control\\
Unexecuted residual & Share of batch volume on the heavier side that
cannot be matched at $\clr$ and remains unexecuted; always rolled over
to the next batch in this implementation, which has no
cancel-on-no-clear path & control\\
\bottomrule
\end{tabularx}
\end{table}

\subsection{Interpretive Status of the RQ2.3 Metrics}
\label{sec:testable}

Three metrics in \Cref{tab:metrics3} require an interpretive caveat that
is easy to state and important not to forget when reading
\Cref{ch:results}.

Order size inflation, boundary concentration, and to a lesser extent
trader surplus are all quantities whose theoretical interest lies in
their being \emph{responses} to the mechanism. Order size inflation is
the signature the literature associates with \emph{pro rata} allocation
specifically; the rule actually implemented here is
price-time-within-the-batch (\Cref{sec:environment}), under which the
mechanically relevant response is timing relative to the batch boundary,
not size, so this metric is reported as a descriptive baseline rather
than as evidence for or against the size-inflation channel. Boundary
concentration is the metric that channel would predict for this rule:
participants clustering submissions near a predictable gate to win
time-priority within the batch.

The orders replayed here were submitted under Hyperliquid's own
continuous engine. Their sizes and arrival times were not chosen in
response to a batch boundary or a batch rationing rule that did not
exist when they were submitted. What the simulation measures is
therefore the mechanical effect of the rule actually implemented on
fixed historical flow: how price-time-within-the-batch would have
allocated these particular orders, and how this particular arrival-time
distribution happens to fall relative to a boundary chosen for the
simulation. These are meaningful quantities --- they establish the
baseline against which a behavioural response would have to be measured
--- but they are not evidence about how participants would behave.
Predictions P5, P6, and P7 of \Cref{tab:predictions} are therefore out
of reach of this design, and \Cref{sec:futurework} discusses the
agent-based extension that would reach them.

The throughput and clearing latency metrics serve a different purpose.
They are not results about market quality but a control: if the two
engines differed substantially in processing performance, a measured
difference in market outcomes could reflect implementation rather than
mechanism. Reporting them is how that alternative explanation is ruled
out.

\section{Validation of the Harness}
\label{sec:validation}

A simulation result is only as credible as the machinery that produced
it. Two layers of validation are applied; two further checks that would
strengthen the case for the control condition were considered but not
built, and are named honestly as gaps rather than as completed work.

\paragraph{Structural invariants, checked live.} The book implementation
enforces a small number of data-structure invariants continuously while
it runs --- for instance, that a price level indexed in the book is
never left empty --- any violation panics the process immediately rather
than allowing a quietly wrong log to be written. The financial
properties a reader might expect here instead --- conservation of
matched quantity, no execution outside a participant's limit, no
duplicated fill --- are not asserted live on every run; they are
verified once, ahead of any run, by the test suite described next, and
the fact that the run behind \Cref{ch:results} completed without the
process aborting (\Cref{tab:validation}) is consistent with, though not
a re-verification of, those properties holding throughout it.

\paragraph{Tests against hand-computed cases.} Two independent
verification passes cover the engines: the codebase's inline test suite
(quantity conservation, price bounds, and correct rationing among them,
across matching scenarios and randomized differential checks against a
naive reference implementation), and a second, self-contained checklist
compiled into the binary itself, that replays small books whose clearing
price and allocation were worked out by hand, including the awkward
cases --- an empty book, no crossing orders, a clearing interval
requiring each stage of the tie-break cascade, and a batch in which a
single order is pivotal. \Cref{tab:validation} reports both pass counts
for the exact build used to produce \Cref{ch:results}' data.

\paragraph{Not built: replay fidelity against the venue.} Comparing the
continuous engine's output trade-for-trade against the trades
Hyperliquid's own engine actually produced from the same flow would be
informative evidence that the control condition is a faithful stand-in
for the mechanism that generated the data. No such comparison pipeline
exists in this codebase; \Cref{tab:validation} marks it not executed
rather than reporting an agreement rate that was never computed.

\paragraph{Not built: degenerate-parameter check.} As $\tau \to 0$ the
FBA should converge toward the CDA in cleared volume and price path,
since a batch containing at most one event is a continuous match.
Verifying this numerically would be a strong end-to-end check of the
batch engine, but it requires the interval sweep that
\Cref{sec:experiments} states was not run; it is listed in
\Cref{sec:futurework} alongside the sweep itself.
